\documentclass[10pt,a4paper]{amsart}
\usepackage{amsmath,amssymb,amsthm,graphicx}
\usepackage{amsmath}
\usepackage{amssymb}
\usepackage{hyperref}
\newtheorem{Theorem}{Theorem}[section]
\newtheorem{Corollary}[Theorem]{Corollary}
\newtheorem{Lemma}[Theorem]{Lemma}
\newtheorem{Proposition}[Theorem]{Proposition}
{\theoremstyle{definition}
\newtheorem{Note}[Theorem]{Note}
\newtheorem{Definition}[Theorem]{Definition}
\newtheorem{Problem}[Theorem]{Problem}
}
\title{Explicit nonzero elements of the algebra O and a nonzero intersection with J}
\author{Paul Terwilliger}
\date{doi:10.3842/SIGMA.2011.069; CC BY4.0}
\begin{document}
\maketitle
\noindent\textit{Editorial note.} All delivered mathematical text is the author's own, copied from the pinned publisher source sources/sigma-2011-069-askey-wilson/pinned-original-source.tex; no statement or proof step was rewritten, reordered or removed.\par\medskip
\noindent\textit{Editorial note.} Preamble provenance: the retained preamble is the source's own layout adapter plus the author's macro block at lines 10-18 of sources/sigma-2011-069-askey-wilson/pinned-original-source.tex (everything between the publisher class metadata and \textbackslash{}begin\{document\}, comments excluded). Nothing was added, no macro was invented, and the publisher document class is deliberately not reproduced; the excerpt serves only as preamble, never as a span coordinate.\par\medskip
\noindent\textit{Editorial note.} External references are navigation pointers only; every deliberately kept step of the author's argument stays inside the counted spans. The displayed numbers of those pointers are the author's own numbering under the author's \textbackslash{}newtheorem declarations (the source has no subsections, so numbering is section.index), computed mechanically by the lane's scanner.\par\medskip
\noindent\textit{Editorial note.} The counted proof uses one bibliography citation (Vidar, Theorem 9.1) and no reference to any labelled item outside its own range, so it needs no navigation stub; the cited entry is reproduced verbatim in this card's own bibliography.tex excerpt, because the bibliography staged in the source directory (faithfulness-bibliography.tex) holds a single entry and does not contain it.\par\medskip
\noindent\textit{Editorial note.} The final sentence of the proof closes the second clause of the statement (that O cap J is nonzero); the author writes it as following from the two computed nonvanishing statements, and it is kept verbatim.\par\medskip
\noindent\textit{Editorial note.} Licence binding: the card's source\_version now carries the value recorded in this source's own licence evidence (sources/sigma-2011-069-askey-wilson/license-observation.json), namely 'SIGMA7(2011),069; publisher final form2011-07-09'. The adopt step of the publication line compares the two and refused the earlier revision; the card is aligned to the evidence rather than the evidence to the card.\par\medskip
\noindent\textit{Allowed standard prerequisites (editorial).}\begin{itemize}
\item The algebra O generated by X, Y and the idempotents xi\_1, xi\_2 (as set up earlier in the source). Used throughout the counted statement and proof. Fixes the elements whose nonvanishing is asserted and the module category in which the verification is made.
\item Vidar's four-dimensional tridiagonal-pair module. Cited inside the counted proof at lines 2572-2577. Supplies the explicit four-dimensional O-module, with the matrices of X and Y displayed in the span, on which the commutators are tested.
\end{itemize}\medskip
\begin{Lemma}
\label{lem:nonz}
The following  elements of~$\mathcal O$ are nonzero:
\begin{gather*}
\xi_1 X \xi_2-\xi_2 X \xi_1,
\qquad
\xi_1 Y \xi_2-\xi_2 Y \xi_1.
\end{gather*}
Moreover
${\mathcal O} \cap J$ is nonzero.
\end{Lemma}

\begin{proof}
To show $\xi_1 X \xi_2-\xi_2 X \xi_1$ is nonzero
we display an $\mathcal O$-module
on which
$\xi_1 X \xi_2-\xi_2 X \xi_1$ is nonzero.
This $\mathcal O$-module is a variation on an $\mathcal O$-module
due to M.~Vidar
\cite[Theorem 9.1]{Vidar}.
Def\/ine $\theta_i = q^{2i}+q^{-2i}$
for $i=0,1,2$. Def\/ine $\vartheta = (q^4-q^{-4})(q^2-q^{-2})(q-q^{-1})^2$.
Adapting~\cite[Theorem 9.1]{Vidar}
there exists a four-dimensional $\mathcal O$-module $V$
with the following property: $V$ has a~basis with respect to
which the matrices representing~$X$, $Y$ are
\begin{gather*}
X:\;
\left(
\begin{array}{ c c c c}
\theta_0 & 0 & 0 & 0   \\
1 & \theta_1 & 0 & 0  \\
0 & 0 & \theta_1  & 0 \\
0 & 1 & 0  & \theta_2
\end{array}
\right),
\qquad
Y:\;
\left(
\begin{array}{ c c c c}
\theta_0 & \vartheta & q & 0   \\
0 & \theta_1 & 0 & 0  \\
0 & 0 & \theta_1  & 1 \\
0 & 0 & 0  & \theta_2
\end{array}
\right).
\end{gather*}
Consider the matrix that represents
$\xi_1X\xi_2-\xi_2 X \xi_1 $ with respect to the above basis.
For this matrix the $(4,3)$-entry is $-q^2$.
This entry is nonzero
so
$\xi_1 X \xi_2- \xi_2 X \xi_1$ is nonzero.
Interchanging the roles of $X$, $Y$ in the above argument,
we see that
 $\xi_1 Y \xi_2- \xi_2 Y \xi_1$ is
nonzero.
The result follows.
\end{proof}

% Bibliography excerpt for this card, copied verbatim from the pinned publisher source
% sources/sigma-2011-069-askey-wilson/pinned-original-source.tex, whose own
% \begin{thebibliography} block starts at line 3560. Only the entry cited inside this
% card's counted span is reproduced, character for character, with no reformatting;
% the wrapper class supplies hyperref, which the entry uses.
\begin{thebibliography}{99}

\bibitem{Vidar}
Vidar M.,
Tridiagonal pairs of shape $(1,2,1)$,
\href{http://dx.doi.org/10.1016/j.laa.2008.02.034}{{\it Linear Algebra Appl.}} {\bf 429} (2008), 403--428,
\href{http://arxiv.org/abs/0802.3165}{arXiv:0802.3165}.

\end{thebibliography}
\end{document}
